\documentclass[11pt]{article}
\usepackage[T1]{fontenc}
\usepackage{lmodern}
\usepackage[letterpaper,margin=0.95in]{geometry}
\usepackage{amsmath,amssymb,amsthm,mathtools}
\usepackage{booktabs}
\usepackage[hidelinks]{hyperref}
\newtheorem{theorem}{Theorem}
\newtheorem{proposition}[theorem]{Proposition}
\newtheorem{corollary}[theorem]{Corollary}
\theoremstyle{remark}
\newtheorem*{remark}{Remark}
\newcommand{\R}{\mathbb R}
\newcommand{\Z}{\mathbb Z}
\newcommand{\sech}{\operatorname{sech}}
\newcommand{\csch}{\operatorname{csch}}
\newcommand{\sinc}{\operatorname{sinc}}
\newcommand{\diag}{\operatorname{diag}}
\newcommand{\norm}[1]{\left\lVert #1\right\rVert}
\title{Neighboring readouts, conditioning, and the scope of a slope barrier}
\author{}
\date{Section 3 source note 5 --- 25 September 2026}

\usepackage{microtype,fancyhdr}
\pagestyle{fancy}\fancyhf{}
\fancyhead[L]{\small Neighboring readouts and conditioning}
\fancyhead[R]{\small Section 3 source note 5}
\fancyfoot[C]{\thepage}
\setlength{\headheight}{14pt}\setlength{\emergencystretch}{2em}
\allowdisplaybreaks[2]
\begin{document}
\maketitle

\begin{abstract}
For an equal-slope tanh dictionary on a uniform lattice, neighboring
differences have whole-line Riesz bounds independent of the number of
features after a common spatial scale is removed.  Physical tanh
readouts on the same zero-sum function space incur an additional
condition-number factor proportional to width.  This improvement leaves
exponential small-slope attenuation intact.  We prove these statements,
identify a distinct bias--halo obstruction on a finite observation
window, and record the exact gradient metric induced by invertible
readout coordinates.  These results supply controls for interpreting a
slope-dependent optimization barrier; they do not establish a universal
training lower bound or an error floor.
\end{abstract}

\section{Dictionary, norm, and the comparison being made}

Fix a center spacing $h>0$ and a dimensionless slope $\lambda>0$.
The physical slope is $\gamma=\lambda/h$.  For $j=0,\ldots,m$, define
\[
 \phi_j(x)=\tanh\!\left(\lambda\left(\frac{x}{h}-j\right)\right),
 \qquad \psi_j=\phi_j-\phi_{j+1}\quad(0\le j<m).
\]
The neighboring-readout map is
\[
 B_m:\R^m\longrightarrow L^2(\R,dx),
 \qquad B_mq=\sum_{j=0}^{m-1}q_j\psi_j.
\]
Coefficients have the Euclidean norm.  Singular values refer to this
coefficient norm and the stated function norm; their ratio is denoted
$\kappa$.  Thus a Gram matrix has condition number $\kappa^2$.
Using $dx/2$ instead of $dx$ multiplies all squared singular values by
$1/2$ and leaves $\kappa$ unchanged.

An individual $\phi_j$ is not in $L^2(\R)$, because it has nonzero tails.
The physical readout comparison must therefore use
\[
 \mathcal W_m=\left\{w\in\R^{m+1}:\sum_{j=0}^m w_j=0\right\},
 \qquad P_mw=\sum_{j=0}^{m}w_j\phi_j.
\]
Every such combination has exponentially decaying tails.  The maps
$P_m$ and $B_m$ represent exactly the same functions, with different
coefficient norms.  Neither includes a bias or an unconstrained final
tanh anchor.  Those columns require a finite observation domain and a
separate analysis below.

\section{A width-independent whole-line bound}

\begin{theorem}[Explicit Riesz bounds]\label{thm:riesz}
For every $m\ge1$, $h>0$, $\lambda>0$, and $q\in\R^m$,
\begin{equation}\label{eq:riesz}
 h a_\lambda\norm{q}_2^2
 \le \norm{B_mq}_{L^2(\R)}^2
 \le 4h\norm{q}_2^2,
 \qquad
 a_\lambda=\frac{4\pi^2}
 {\lambda^2\sinh^2(\pi^2/(2\lambda))}.
\end{equation}
Consequently,
\begin{equation}\label{eq:C}
 \sqrt{h a_\lambda}\le \sigma_{\min}(B_m)
 \le \sigma_{\max}(B_m)\le2\sqrt h,
 \qquad
 \kappa(B_m)\le C_\lambda
 :=\frac{\lambda}{\pi}\sinh\!\left(\frac{\pi^2}{2\lambda}\right).
\end{equation}
In particular, the bounds for $B_m/\sqrt h$ do not depend on width.
\end{theorem}

\begin{proof}
Use the dimensionless coordinate $t=x/h$.  Introduce the probability
density and the piecewise constant coefficient function
\[
 \rho_\lambda(t)=\frac{\lambda}{2}\sech^2(\lambda t),
 \qquad p_q(t)=\sum_{j=0}^{m-1}q_j\mathbf 1_{[j,j+1]}(t).
\]
Integration of the tanh derivative gives the exact identity
\begin{equation}\label{eq:convolution}
 (B_mq)(ht)=2(\rho_\lambda*p_q)(t).
\end{equation}
Since $\norm{\rho_\lambda}_1=1$ and
$\norm{p_q}_2^2=\norm{q}_2^2$, the $L^1$--$L^2$ convolution inequality
proves the upper bound in \eqref{eq:riesz}.

Use the Fourier convention
$\widehat f(\omega)=\int_\R f(t)e^{-i\omega t}\,dt$.
The transform of the smoothing density is
\begin{equation}\label{eq:multiplier}
 \widehat\rho_\lambda(\omega)=M_\lambda(\omega)
 =\frac{\pi\omega/(2\lambda)}{\sinh(\pi\omega/(2\lambda))},
 \qquad M_\lambda(0)=1.
\end{equation}
For completeness, setting $u=e^{2\lambda t}$ and
$a=\omega/(2\lambda)$ reduces the transform to
\[
 \int_0^\infty\frac{u^{-ia}}{(1+u)^2}\,du
 =\Gamma(1-ia)\Gamma(1+ia)
 =\frac{\pi a}{\sinh(\pi a)}.
\]
The first equality is the beta integral and the second is the gamma
reflection identity, continued at $a=0$.

Writing $\sinc u=\sin(u)/u$, with $\sinc0=1$, we also have
\[
 \widehat p_q(\omega)=e^{-i\omega/2}\sinc(\omega/2)Q(\omega),
 \qquad Q(\omega)=\sum_{j=0}^{m-1}q_je^{-ij\omega}.
\]
For $|\omega|\le\pi$, $\sinc(\omega/2)\ge2/\pi$ and
$M_\lambda(\omega)\ge M_\lambda(\pi)$: the function $v/\sinh v$
decreases for $v\ge0$.  Plancherel's identity, restriction to this
frequency interval, and Parseval's identity for $Q$ give
\begin{align*}
 \frac1h\norm{B_mq}_2^2
 &=\frac4{2\pi}\int_\R
 M_\lambda(\omega)^2\sinc^2(\omega/2)|Q(\omega)|^2\,d\omega\\
 &\ge\frac{16}{\pi^2}M_\lambda(\pi)^2
 \frac1{2\pi}\int_{-\pi}^{\pi}|Q(\omega)|^2\,d\omega
 =a_\lambda\norm{q}_2^2.
\end{align*}
Taking square roots and ratios proves \eqref{eq:C}.
\end{proof}

The word ``Riesz'' means precisely the two norm bounds
\eqref{eq:riesz}.  They hold for every finite coefficient sequence,
including sequences placed at any integer translate of the displayed
lattice segment.  Absolute singular values still scale as $\sqrt h$;
the theorem controls their ratio, rather than asserting a nonvanishing
absolute scale under grid refinement.

\section{The physical readout penalty and the sharp-step limit}

Define the rectangular difference map $E_m:\R^m\to\mathcal W_m$ by
\[
 E_mq=(q_0,\ q_1-q_0,\ldots,q_{m-1}-q_{m-2},\ -q_{m-1})^T.
\]
It is a bijection onto $\mathcal W_m$, and
$B_m=P_mE_m$.  Both endpoint terms matter.

\begin{proposition}\label{prop:physical}
The singular values of $E_m$, in increasing order, are
\begin{equation}\label{eq:Esingular}
 e_k=2\sin\!\left(\frac{k\pi}{2(m+1)}\right),
 \quad 1\le k\le m,
 \qquad
 \kappa(E_m)=\cot\!\left(\frac{\pi}{2(m+1)}\right).
\end{equation}
Moreover,
\begin{equation}\label{eq:physical}
 \frac{\kappa(E_m)}{C_\lambda}
 \le\kappa(P_m)\le C_\lambda\kappa(E_m).
\end{equation}
Thus $\kappa(P_m)=\Theta_\lambda(m)$ at fixed $\lambda>0$, whereas
$\kappa(B_m)$ stays bounded as $m$ increases.
\end{proposition}

\begin{proof}
The matrix $E_m^TE_m$ has diagonal entries $2$ and off-diagonal entries
$-1$.  Its sine eigenvectors, with entries
$\sin(jk\pi/(m+1))$ for $1\le j\le m$, have eigenvalues
$4\sin^2(k\pi/(2(m+1)))$.  This proves \eqref{eq:Esingular}.
For $w=E_mq\ne0$, Theorem~\ref{thm:riesz} gives
\[
 h a_\lambda\frac{\norm q_2^2}{\norm{E_mq}_2^2}
 \le\frac{\norm{P_mw}_2^2}{\norm w_2^2}
 \le4h\frac{\norm q_2^2}{\norm{E_mq}_2^2}.
\]
Taking maxima and minima over $q$ yields
\[
 \frac{\sqrt{h a_\lambda}}{e_1}\le\sigma_{\max}(P_m)
 \le\frac{2\sqrt h}{e_1},
 \qquad
 \frac{\sqrt{h a_\lambda}}{e_m}\le\sigma_{\min}(P_m)
 \le\frac{2\sqrt h}{e_m}.
\]
Dividing the respective bounds proves \eqref{eq:physical}.
\end{proof}

This comparison identifies a cumulative-coordinate penalty.  It does
not assert that every singular value increases under differencing, or
that $\kappa(B_m)<\kappa(P_m)$ for every finite pair $(m,\lambda)$.

There is an exact limiting comparison.  For fixed $m,h$, as
$\lambda\to\infty$, the densities $\rho_\lambda$ form an approximate
identity.  Equation~\eqref{eq:convolution} therefore converges in $L^2$
to $2p_q$, equivalently
$\psi_j\to2\mathbf1_{[jh,(j+1)h]}$ in $L^2(dx)$.  In the resulting
sharp-step dictionary,
\begin{equation}\label{eq:step}
 B_m^*B_m=4hI,
 \qquad \kappa(B_m)=1,
 \qquad \kappa(P_m)=\kappa(E_m)
 \sim\frac{2(m+1)}{\pi}.
\end{equation}
The last identity follows from
$\norm{P_mE_mq}_2^2=4h\norm q_2^2$: the singular values of $P_m$
are $2\sqrt h/e_k$, in reverse order.  The Gram condition number
changes from order $m^2$ to one.  The bound $C_\lambda$ is conservative
and tends to $\pi/2$, so this sharp limit uses the exact argument above.

\section{Width uniformity does not remove small-slope attenuation}

To see which difficulty remains, periodize the Fourier energy.
For every finite $q$,
\begin{equation}\label{eq:gramian}
 \norm{B_mq}_2^2=\frac h{2\pi}\int_{-\pi}^{\pi}
 G_\lambda(\omega)|Q(\omega)|^2\,d\omega,
 \quad
 G_\lambda(\omega)=4\sum_{k\in\Z}
 \sinc^2\!\left(\frac{\omega+2\pi k}{2}\right)
 M_\lambda(\omega+2\pi k)^2.
\end{equation}
This continuous, positive multiplier describes the infinite translated
dictionary.  Theorem~\ref{thm:riesz} gives
$a_\lambda\le G_\lambda\le4$.  At the constant and alternating
coefficient frequencies,
\begin{equation}\label{eq:twofreq}
 G_\lambda(0)=4,
 \qquad
 G_\lambda(\pi)=\frac{8\pi^2}{\lambda^2}
 \sum_{r=0}^{\infty}
 \csch^2\!\left((2r+1)\frac{\pi^2}{2\lambda}\right).
\end{equation}
Indeed, all nonzero aliases vanish at zero frequency.  At $\pi$,
$\sinc^2(\nu/2)=4/\nu^2$ for odd multiples $\nu$ of $\pi$, which
cancels the numerator of $M_\lambda(\nu)^2$; pairing positive and
negative frequencies gives the series.

Define the infinite-lattice condition number by
\[
 \kappa_\infty(\lambda)
 =\sqrt{\frac{\max_{[-\pi,\pi]}G_\lambda}
 {\min_{[-\pi,\pi]}G_\lambda}}.
\]
No assertion about the location of the minimum is needed for
\begin{equation}\label{eq:infinitebounds}
 \sqrt{\frac4{G_\lambda(\pi)}}
 \le\kappa_\infty(\lambda)\le C_\lambda.
\end{equation}
These modes also have a finite-sequence meaning.  The real packets
$q_j=m^{-1/2}$ and $q_j=(-1)^jm^{-1/2}$ have Fourier energy tending
to point masses at $0$ and $\pi$, respectively.  This is the
approximate-identity property of the Fej\'er kernel $|Q|^2$.
Their squared gains tend to $hG_\lambda(0)$ and $hG_\lambda(\pi)$.
In particular,
$\liminf_{m\to\infty}\kappa(B_m)\ge\sqrt{4/G_\lambda(\pi)}$.
An infinite constant sequence itself is not an $\ell^2$ coefficient
vector and is not being used as one.

Set $A=\pi^2/(2\lambda)$.  As $\lambda\downarrow0$, the first term in
\eqref{eq:twofreq} dominates, since
$\csch^2((2r+1)A)=4e^{-2(2r+1)A}(1+o(1))$.  Hence
\begin{equation}\label{eq:smalllambda}
 G_\lambda(\pi)\sim\frac{32\pi^2}{\lambda^2}e^{-\pi^2/\lambda},
 \quad
 \sqrt{\frac4{G_\lambda(\pi)}}
 \sim\frac{\lambda}{2\sqrt2\,\pi}e^{\pi^2/(2\lambda)},
 \quad
 C_\lambda\sim\frac{\lambda}{2\pi}e^{\pi^2/(2\lambda)}.
\end{equation}
Thus the remaining exponential dependence on small $\lambda$ is real,
not merely an artifact of the lower-bound proof.  For example,
$\lambda=0.25$ gives $a_\lambda\approx1.80834\times10^{-14}$,
$C_\lambda\approx1.48727\times10^7$, and the lower bound in
\eqref{eq:infinitebounds} is approximately $1.05166\times10^7$.
These are whole-line spectral constants, not condition numbers of a
particular finite training matrix.  The large-width packet limit and
the small-$\lambda$ limit are distinct statements.

\section{Finite windows, halo columns, and the missing lower bound}

Restricting an integral to $[-1,1]$ decreases its value.  Therefore
\eqref{eq:riesz} supplies a finite-window upper bound but does not
supply the same lower bound.  A bump supported mostly outside the
window can have arbitrarily little observed energy.  Bias and anchor
columns introduce an additional concrete obstruction.

\begin{proposition}[Bias--halo cancellation]\label{prop:halo}
Let $\mathsf B$ be a real sampled feature matrix with $M$ rows, at
least as many rows as columns, and samples $x_i\in[-1,1]$.  Suppose it
contains the columns
\[
 \frac{d_b}{\sqrt M}\mathbf1,
 \qquad \frac{t}{\sqrt M}\phi_W(x_i),
 \qquad
 \phi_W(x)=\tanh\!\left(\frac\lambda h(x-x_W)\right),
 \quad x_W=1+Rh,
\]
where $d_b,t,\lambda,h>0$ and $R\ge0$.  There is no assumption on its
other columns.  Then
\begin{equation}\label{eq:halo}
 \sigma_{\min}(\mathsf B)
 \le\frac{2t}{\sqrt{1+(t/d_b)^2}}e^{-2\lambda R},
 \qquad
 \kappa(\mathsf B)
 \ge\frac{\sqrt{d_b^2+t^2}}{2t}e^{2\lambda R}.
\end{equation}
The full condition number is infinite if the matrix is rank deficient.
\end{proposition}

\begin{proof}
For every observed $x$,
\[
 0<1+\phi_W(x)
 =\frac{2e^{2\lambda(x-x_W)/h}}{1+e^{2\lambda(x-x_W)/h}}
 \le2e^{-2\lambda R}.
\]
Choose a coefficient vector $v$ equal to $t/d_b$ in the bias position,
$1$ in the anchor position, and zero elsewhere.  Then
$\norm v_2^2=1+(t/d_b)^2$ and
$\norm{\mathsf Bv}_2\le2t e^{-2\lambda R}$.
This bounds $\sigma_{\min}$ from above.  The bias column gives
$\sigma_{\max}\ge d_b$, proving the second inequality.
\end{proof}

In particular, $R=\lceil\sqrt N\rceil$ gives a lower bound at least
$\tfrac12 e^{2\lambda\lceil\sqrt N\rceil}$ for the full condition
number.  Increasing $\lambda$ can improve the interior smoothing
spectrum while worsening this out-of-domain cancellation.  Neither
bound identifies which directions a target residual occupies.

Three further changes require their own estimates.  Scaling the bump
coefficients by a positive diagonal $S$ gives only
$\kappa(B_mS)\le C_\lambda S_{\max}/S_{\min}$; the scale ratio may
depend on width.  Replacing integration by samples requires an
operator-norm bound on the Gram approximation error to preserve a
positive lower bound.  Finally, unequal slopes invalidate the common
convolution multiplier; same-sign slopes still cancel tails, but
opposite signs need not.  None of these cases is covered merely by
citing Theorem~\ref{thm:riesz}.

\section{Invertible coordinates and what they control}

For a finite window, let
$c=(b,w_1,\ldots,w_W)^T$ be the physical coefficients and define the complete sampled dictionary by
\[\mathsf A=M^{-1/2}[\mathbf1,\phi_1,\ldots,\phi_W].\]  Geometry is frozen here.  For any fixed invertible
matrix $T$, set $c=T\theta$.  Then
\begin{equation}\label{eq:metric}
 \operatorname{range}(\mathsf AT)=\operatorname{range}(\mathsf A),
 \qquad
 \nabla_\theta\mathcal L=T^T\nabla_c\mathcal L,
 \qquad
 \Delta c=-\eta TT^T\nabla_c\mathcal L
\end{equation}
for one ordinary gradient-descent step in $\theta$.  The first identity
preserves approximation capacity.  The last identity changes the
physical gradient metric, and is exact for any differentiable loss
with this fixed linear map.  It is not the update rule for Adam.

Here are the concrete maps used in the source implementation review.
Let $\alpha_j>0$ and $\alpha_b>0$ be fixed reference allowances; they
are scales, not constraints imposed on trained coefficients.  Put
$s_j=\sum_{k=1}^j\alpha_k$.  Let $C$ be the square cumulative-sum
matrix and $L=C^{-1}$ its lower bidiagonal inverse.  In each row of
Table~\ref{tab:maps}, $T=\diag(t_b,T_w)$.
\begin{table}[ht]
\centering
\caption{Invertible maps for the complete physical readout.}
\label{tab:maps}
\begin{tabular}{lll}
\toprule
Coordinates & Hidden-readout map $T_w$ & Bias scale $t_b$\\
\midrule
Raw & $I$ & $1$\\
Collective & $\diag(\sqrt{\alpha_j})$ & $\sqrt{\alpha_b}$\\
Individual & $\diag(\alpha_j)$ & $\alpha_b$\\
Collective neighboring & $L\diag(\sqrt{s_j})$ & $\sqrt{\alpha_b}$\\
Individual neighboring & $L\diag(s_j)$ & $\alpha_b$\\
\bottomrule
\end{tabular}
\end{table}

For a neighboring map $w=LS\theta_w$, the feature row is exactly
\[
 [\phi_1,\ldots,\phi_W]LS
 =[S_{11}(\phi_1-\phi_2),\ldots,
 S_{W-1,W-1}(\phi_{W-1}-\phi_W),S_{WW}\phi_W].
\]
The final anchor makes $L$ square and invertible.  This map differs
from the rectangular $E_m$ used for the whole-line zero-sum comparison.
It induces the coupled physical metric $LS^2L^T$ for hidden readouts.
The scales multiply cumulative coefficients before differencing:
unequal scales applied to individual tanhs first would leave uncancelled
tails in $d_j\phi_j-d_{j+1}\phi_{j+1}$.

A matched physical start must be encoded as
$\theta_0=T^{-1}c_0$.  Independent random draws in each coordinate
system generally give different physical functions.  Likewise, native
isotropic damping changes its physical penalty under this map:
$\mu\norm{\Delta\theta}_2^2
=\mu\norm{T^{-1}\Delta c}_2^2$.  A claim about coordinates should
specify both initialization and the optimizer or regularizer.

Let $y\in\R^M$ contain the target values at the samples. For the squared training loss,
$\mathcal L(\theta)=\tfrac12\norm{\mathsf AT\theta-y/\sqrt M}_2^2$,
the Hessian is $(\mathsf AT)^T(\mathsf AT)$.  With
$0<\eta<2/\norm{\mathsf AT}_2^2$, an attainable left-singular residual
component evolves exactly as
\begin{equation}\label{eq:rate}
 e_i(k)=(1-\eta\sigma_i^2)^k e_i(0),
\end{equation}
where $\sigma_i>0$ is the corresponding singular value of
$\mathsf AT$.  Components orthogonal to its range remain unchanged.
A tiny positive singular value is therefore a rate obstruction for
this algorithm when its residual component is nonzero; it is not an
exact-arithmetic error floor.  Stability and rates use the complete
matrix, including bias, anchor, and halos, not the isolated bound
$4h$ from Theorem~\ref{thm:riesz}.

These facts locate the role of neighboring readouts in a proposed
``$\gamma$ barrier.''  At fixed spacing, changing $\gamma$ changes
$\lambda=h\gamma$ and hence the dictionary's smoothing.  Changing
$T$ at fixed $\gamma$ preserves its represented function space but
changes the gradient metric and its singular spectrum.  A frozen
comparison at matched physical initialization can therefore separate
some coordinate dependence from slope dependence.  Reference
allowances should be held fixed during that comparison if the intended
change is slope alone.  The whole-line calculation predicts a
small-$\lambda$ obstruction in fine coefficient modes after removing
the cumulative-coordinate penalty.  It does not prove that a given
target loads those modes, that finite-window halos are harmless, or
that a measured fraction of finite-budget optimization error is caused
by bounded slopes.  Such conclusions need the actual complete
dictionary and residual, with the algorithm specified.

\paragraph{Source and scope.}
This standalone note reorganizes the analytic argument in
\path{docs/neighbor_difference_conditioning.md} and the readout-map
audit in
\path{docs/frozen_geometry_capacity_access/pr_conditioning_review.md}.
The latter records the maps in the merged implementation at commit
\texttt{d7d26e6}.  The numerical constants above are evaluations of the
displayed formulas, as also recorded in the former source.  No new
experiment or empirical convergence claim is included here.

\end{document}
